\chapter{Vor-/Nachteile und Grenzen}
\label{ch:zsfg}
Als Abschluss meiner Seminararbeit möchte ich einmal auf die Vorteile, Nachteile und die Grenzen der Anomalieerkennung in adaptiven Systemen mit Generativen KI zum heutigen Stand eingehen. Zu Beginn werde ich mit den Vorteilen bei der Nutzung von generativer KI in der \\ Anomalieerkennung schreiben. Ein großer Vorteil von generativer KI bei der Anomalieerkennung ist wie in meinem Fallbeispiel davor erläutert, dass es eine große Verbesserung der Genauigkeit bei der Erkennung von neuen oder auch bereits gelernten Anomalien gibt. Der nächste Vorteil baut auf die verbesserte Genauigkeit auf und beinhaltet eine bessere Anpassungsfähigkeit des Systems das eine generative KI benutzt. Durch das Frühzeitige entdecken neuer oder bekannter Anomalien kann ein adaptives System besser und schneller Anpassungsstrategien entwickeln und umsetzen um einen größeren Schaden entweder komplett zu verhindern oder diesen stark zu minimieren. Ein weiterer Vorteil ist die Analyse und die Entdeckung von Anomalien mit Echtzeit Daten. \\ Da Systeme in der heutigen Zeit immer größer und komplexer werden, gerade mit Hinblick auf Big-Data ist es essenziell eine automatisierte und verlässliche Erkennung von Anomalien in \\ Echtzeit zu unterstützen, da die herkömmlichen statischen Methoden an Ihre Grenzen kommen. Die Datenmengen sind in heutigen Systemen so groß, dass es unmöglich ist Anomalien in Echtzeit vorherzusagen und das System korrekt anzupassen. Durch den Einsatz von generativer KI werden auch Realistische Daten generiert die dem System helfen sich besser flächendeckender \\ vor Fehlern die durch Anomalien entstehen zu schützen. Das System ist robuster gegenüber \\ Neuheiten, da diese eventuell durch das generieren von realistischen Daten bereits zu einem Teil abgedeckt wurden \cite{Fallbeispiel}\cite{Roadmap}. 
\nAbsatz
Weitergehend möchte ich auf die Nachteile eingehen. Ein großer Punkt ist die Ethische \\ Herausforderung für die Systeme. Bereits gelernte Muster können dazu tendieren dinge zu \\ generieren die Menschen aus bestimmten Gruppen diskriminieren. Dabei ist es ebenfalls sehr schwer herauszufinden ob ein Modell bereits Vorurteile gebildet hat und wie es Sie verwendet\cite{Bias}. Damit kommen wir auch bereits zum nächsten Punkt. Die heutigen komplexen Generativen \\ Modelle sind nur schwer nachzuvollziehen. Je größer das Modell ist und je mehr Trainingsdaten verwendet wurden desto schwieriger wird es zu verstehen warum welche Entscheidung \\ getroffen wurde. Durch die Fehlende Transparenz ist man meist gezwungen ein Ergebnis zu \\ akzeptieren wenn es den Erwartungen oder tests entspricht\cite{explainability}. Ein weiteres Problem ist die Prompt Robustheit. Damit ist gemeint dass man einer KI immer wieder den selben Prompt als Eingabe \\ geben kann aber dennoch wird das Ergebnis durch das hinzugefügt Rauschen unterschiedlich sein. Damit ist es schwierig Fehler oder allgemein Outputs zu rekonstruieren\cite{robustness}. Ein letzter Nachteil ist eine Inkorrekte Adaption durch fehlerhafte Daten die zu einem falsches Training des Modells führen. Durch fehlerhafte Adaption des Systems kann es zu Leistungsverschlechterungen \\ kommen die wiederum zu mehr Fehlern im System führen\cite{Adaption}. Auch ein Finanzieller Verlust kann dort zustande kommen. Das folgende Szenario ist selbst ausgedacht um diesen Punkt \\ verständlicher zu machen. In einem Online Shop gibt es einen Sale und viele User besuchen gleichzeitig die Seite. Dabei wäre es Fatal wenn das System seine verfügbaren Ressourcen wie die Kapazität von gleichzeitig verbundenen Usern beschränkt, da es denkt dass ein Angriff folgt, anstatt die Kapazität zu erhöhen. Durch das Drosseln der Kapazität würde es zu finanziellen \\ Verlusten kommen da es einen Anomalie erkennt, obwohl es valide ist. 
\nAbsatz
Zum Schluss meiner Arbeit möchte ich einmal auf die meiner Meinung nach \\ Aktuellen Grenzen eingehen. Ein entscheidender Faktor hierbei wäre das Gleichgewicht zwischen Privatsphäre und Effektivität. Vielen Menschen ist Ihre Privatsphäre sehr wichtig weswegen \\ einige Daten nicht vollständig mit ihrer Zugehörigkeit übermittelt werden dürfen. Hierbei wäre es aber für das Training und die Effektivität des Modells in manchen Fällen sehr hilfreich um bessere und robustere Muster bilden zu können. Jedes Modell ist zudem noch Limitiert durch physische Grenzen. Da dinge wie Rechenleistung, Ressourcen wie Strom oder Geld und Speicherkapazität endlich sind werden Modelle meist durch einen oder mehrere dieser Faktoren begrenzt. Ebenfalls ist eine momentane Grenze die Genauigkeit. Kein Modell hat bislang eine Genauigkeit oder werte in allen Vergleichsmetriken von 100\% . Selbst CloudGEN kommt nur mit 99,8\% nah dran aber erreicht nicht die 100\% . Demnach muss man immer vorsichtig sein, da selbst wenn es nur geringe Lücken sind trotzdem diese gibt. Ein weiterer Punkt hierbei ist noch die Skalierbarkeit. \\ Die heutigen Modelle und Systeme sind sehr groß und natürlich zu einem gewissen teil \\ skalierbar, aber wie bereits erwähnt gibt es auch bei dem skalieren eine physische Grenze. \\ Diese ist an verschiedenen Punkten schneller erreicht als bei anderen. So ist es möglich dass \\ eine große Firma keine Probleme mit den Kosten für den Strom haben wird, aber dafür keinen physischen Platz für die Speicherkapazitäten hat. Zum Schluss würde ich Persönlich sagen dass Generative KI bei der Anomalieerkennung eine wichtige Rolle spielt und zukünftig noch \\ wichtiger wird, aber man immer die Grenzen und nachteile im Auge behalten sollte. Vor allem ist nichts ohne Limit, sei es eine Performance oder auch die Ressourcen. Für die Zukunft wäre es wichtig gerade in der heutigen Zeit auf die knappen aber dennoch verfügbaren Ressourcen die wir haben zu achten und diese effektiv einzusetzen. Dazu sollte man ebenfalls auch noch versuchen die Ergebnisse von den Generativen KIs transparenter zu gestalten, so dass man die  \\Entscheidungen besser nachvollziehen kann. Dennoch ist diese Technologie sehr vielversprechend für unseren bestehenden Systeme und auch wenn sie nicht zu 100\% ausgreift ist, \\ kann sie für zuverlässige Einsätze in der Medizin oder ähnlichen wichtigen Branchen eine große Hilfe werden.    
\nAbsatz
Zusammenfassend lässt sich sagen, dass der Einsatz von generativer KI in vielen Bereich vor allem wenn es um Big-Data geht in der Anomalieerkennung sinn macht und die Performance erheblich verbessert. Dennoch gibt es viele Risiken auf die man Achten muss. Meiner Meinung nach ist dieses Gebiet für die Zukunft sehr wichtig und effizient. Natürlich muss noch definiert werden wo die Privatsphäre des einzelenen anfängt und die Effektivität des Trainings der Modelle aufhört, aber das Potential ist sehr groß. In der Zukunft sollte man sich auf jeden Fall mit der Genauigkeit befassen. Auch wenn es fast unmöglich ist eine 100\% Genauigkeit zu erzielen, kann es trotzdem bei zum Beispiel autonom fahrenden Autos bei einer genauigkeit von 99,9\% und 40 Millionen Autos auf der Straße täglich zu 40.000 Unfällen kommen. Man sollte die Werte mit Vorsicht betrachten und auch wenn sie in einem kleinen Maße sehr vielversprechend klingen sind sie bei einer höheren Skalierung und einem anderen Kontext sehr Kritisch. Ebenfalls sollte man die Erklärbarkeit ein wenig transparenter machen. Durch mehr Transparenz müsste man nicht mehr blind auf die Ergebnisse vertrauen und könnte den Output von Ergebnissen oder auch die Erkennung von Anomalien begründen. Ich denke dass die Benutzung von generativer KI in den nächsten Jahren immer mehr Zuwachs erfahren wird und man sie in unserem Alltag öfter sehen oder mitbekommen wird.    